\chapter{Cyclic Groups and Their Automorphism Groups}

\section{Basic Facts about Cyclic Groups}

\begin{definition}[Cyclic group]
Let $G$ be a group. We say that $G$ is \textbf{cyclic} if there exists an element
$g\in G$ such that
\[
    G=\gen{g}=\{g^k:k\in\Z\}.
\]
If $G$ has order $n$, then it is usually denoted by $C_n$.

A finite cyclic group is often written additively:
\[
    C_n\cong \Z/n\Z=\{\overline{0},\overline{1},\ldots,\overline{n-1}\}.
\]
Since cyclic groups are abelian, both multiplicative and additive notation are commonly used.
\end{definition}

\begin{proposition}[Basic properties of cyclic groups]
Let $G$ be a cyclic group.
\begin{enumerate}[label=(\arabic*)]
    \item A finite cyclic group is uniquely determined up to isomorphism by its order.
    \item Every subgroup of a cyclic group is cyclic.
    \item Every quotient group of a cyclic group is cyclic.
\end{enumerate}
\end{proposition}

\begin{proof}
Suppose $G=\gen{g}$ and $|G|=n$. Then the map
\[
    \gen{g}\longrightarrow \Z/n\Z,\qquad g^i\longmapsto \overline{i}
\]
is an isomorphism. Hence every cyclic group of order $n$ is isomorphic to $\Z/n\Z$.

Now let $H\leq \gen{g}$. If $H=\{1\}$, then $H$ is cyclic. Otherwise, choose the smallest positive integer $d$ such that $g^d\in H$. We claim that
\[
    H=\gen{g^d}.
\]
Indeed, if $g^m\in H$, write $m=qd+r$ with $0\leq r<d$. Then
\[
    g^r=g^m(g^d)^{-q}\in H.
\]
By the minimality of $d$, we must have $r=0$. Hence $g^m\in\gen{g^d}$, and so $H=\gen{g^d}$.

Finally, if $N\lhd G=\gen{g}$, then
\[
    G/N=\gen{gN},
\]
so $G/N$ is cyclic.
\end{proof}

\section{Automorphisms of Finite Cyclic Groups}

\begin{definition}[The unit group modulo $n$]
The group of units modulo $n$ is
\[
    \Unit(n)=(\Z/n\Z)^\times
    =\{\overline{a}\in \Z/n\Z:\gcd(a,n)=1\}.
\]
Its group operation is multiplication modulo $n$, and
\[
    |\Unit(n)|=\varphi(n),
\]
where $\varphi$ is Euler's phi function.
\end{definition}

\begin{theorem}[Automorphisms of a finite cyclic group]\label{thm:aut-cyclic}
Let $C_n=\gen{g}$. Then
\[
    \Aut(C_n)\cong (\Z/n\Z)^\times=\Unit(n).
\]
In particular,
\[
    |\Aut(C_n)|=\varphi(n).
\]
\end{theorem}

\begin{proof}
Every automorphism of $C_n$ is completely determined by the image of the generator $g$.

Let $\sigma\in\Aut(C_n)$. Since $C_n=\gen{g}$, there exists an integer $i$ such that
\[
    \sigma(g)=g^i.
\]
Because automorphisms preserve orders, $\sigma(g)=g^i$ must also have order $n$. This happens exactly when
\[
    \gcd(i,n)=1.
\]
Thus each automorphism determines a unit $\overline{i}\in(\Z/n\Z)^\times$.

Conversely, if $\gcd(a,n)=1$, then $g^a$ is also a generator of $C_n$. Hence the assignment
\[
    g\longmapsto g^a
\]
extends uniquely to an automorphism of $C_n$. Denote this automorphism by $\sigma_a$.

Now define
\[
    \Psi:\Aut(C_n)\longrightarrow (\Z/n\Z)^\times,
    \qquad
    \sigma_a\longmapsto \overline{a}.
\]
This map is bijective by the previous discussion. Moreover,
\[
    (\sigma_a\circ\sigma_b)(g)
    =\sigma_a(g^b)
    =(g^a)^b
    =g^{ab},
\]
so
\[
    \sigma_a\circ\sigma_b=\sigma_{ab}.
\]
Therefore $\Psi$ is a group isomorphism.
\end{proof}

\begin{remark}
Theorem \ref{thm:aut-cyclic} reduces the computation of $\Aut(C_n)$ to the computation of the unit group $\Unit(n)$.
\end{remark}

\section{Odd Prime Powers}

\subsection{The prime case}

\begin{remark}[Finite fields]
Let $\F_{p^e}$ be the finite field with $p^e$ elements. Then
\[
    \F_{p^e}^{+}\cong \underbrace{C_p\oplus\cdots\oplus C_p}_{e\text{ copies}},
\]
and
\[
    \F_{p^e}^{\times}\cong C_{p^e-1}.
\]
In particular,
\[
    \F_p^\times\cong C_{p-1}.
\]
\end{remark}

\begin{remark}
        Let \(p\) be a prime and let \(e\geq 1\). Then
        \[
            \F_{p^e}^{\times}\cong C_{p^e-1}.
        \]
    \end{remark}
        
        \begin{proof}
        Put \(q=p^e\) and \(F=\F_q\). Let
        \[
            G=F^\times.
        \]
        Then \(G\) is a finite abelian group under multiplication, and
        \[
            |G|=q-1=p^e-1.
        \]
        
        Let \(m\) be the exponent of \(G\), namely
        \[
            m=\operatorname{lcm}\{\operatorname{ord}(a):a\in G\}.
        \]
        By Lagrange's theorem, \(\operatorname{ord}(a)\mid |G|=q-1\) for every \(a\in G\). Hence
        \[
            m\mid q-1.
        \]
        In particular,
        \[
            m\leq q-1.
        \]
        
        We first show that there exists an element of \(G\) whose order is exactly \(m\).
        Write
        \[
            m=\prod_{\ell\mid m}\ell^{a_\ell}
        \]
        as a product of prime powers. For each prime \(\ell\mid m\), by the definition of \(m\), there exists \(g_\ell\in G\) such that the \(\ell\)-part of \(\operatorname{ord}(g_\ell)\) is \(\ell^{a_\ell}\). Define
        \[
            h_\ell
            =
            g_\ell^{\operatorname{ord}(g_\ell)/\ell^{a_\ell}}.
        \]
        Then
        \[
            \operatorname{ord}(h_\ell)=\ell^{a_\ell}.
        \]
        Since \(G\) is abelian, the elements \(h_\ell\) commute. Therefore the product
        \[
            g=\prod_{\ell\mid m}h_\ell
        \]
        has order
        \[
            \operatorname{ord}(g)
            =
            \prod_{\ell\mid m}\ell^{a_\ell}
            =
            m,
        \]
        because the orders of the \(h_\ell\)'s are pairwise coprime.
        
        Now, by the definition of \(m\), every element \(a\in G\) satisfies
        \[
            a^m=1.
        \]
        Thus every element of \(G=F^\times\) is a root of the polynomial
        \[
            X^m-1\in F[X].
        \]
        But a nonzero polynomial of degree \(m\) over a field has at most \(m\) roots. Hence
        \[
            q-1=|F^\times|\leq m.
        \]
        Together with \(m\leq q-1\), we get
        \[
            m=q-1.
        \]
        
        Therefore \(G\) contains an element \(g\) of order \(q-1=|G|\). Hence
        \[
            G=\langle g\rangle
        \]
        is cyclic. Consequently,
        \[
            \F_{p^e}^{\times}
            =
            F^\times
            \cong C_{q-1}
            =
            C_{p^e-1}.
        \]
        \end{proof}


\begin{corollary}
If $p$ is prime, then
\[
    \Aut(C_p)\cong \Unit(p)\cong \F_p^\times\cong C_{p-1}.
\]
\end{corollary}

\begin{proof}
By Theorem \ref{thm:aut-cyclic},
\[
    \Aut(C_p)\cong(\Z/p\Z)^\times.
\]
But $(\Z/p\Z)^\times=\F_p^\times$, and the multiplicative group of a finite field is cyclic. Hence
\[
    \Aut(C_p)\cong C_{p-1}.
\]
\end{proof}

\subsection{Odd prime powers}

\begin{theorem}[Units modulo odd prime powers]\label{thm:odd-prime-power}
Let $p$ be an odd prime and let $e\geq 2$. Then
\[
    \Unit(p^e)\cong C_{p^{e-1}(p-1)}.
\]
Equivalently,
\[
    \Aut(C_{p^e})\cong C_{p^{e-1}(p-1)}.
\]
\end{theorem}

\begin{proof}
We show that $\Unit(p^e)$ contains an element of order
\[
    \varphi(p^e)=p^{e-1}(p-1).
\]

First choose an integer $\lambda$ whose image modulo $p$ is a primitive root in $\F_p^\times$. Thus $\lambda$ has order $p-1$ modulo $p$.

Consider $\lambda$ modulo $p^2$. Since its image modulo $p$ has order $p-1$, its order modulo $p^2$ is either $p-1$ or $p(p-1)$.

If
\[
    \lambda^{p-1}\not\equiv 1\pmod{p^2},
\]
then $\lambda$ has order $p(p-1)$ modulo $p^2$.

If instead
\[
    \lambda^{p-1}\equiv 1\pmod{p^2},
\]
then consider $\lambda+p$. By the binomial theorem,
\[
    (\lambda+p)^{p-1}
    \equiv
    \lambda^{p-1}+p(p-1)\lambda^{p-2}
    \pmod{p^2}.
\]
Using $\lambda^{p-1}\equiv1\pmod{p^2}$, we get
\[
    (\lambda+p)^{p-1}
    \equiv
    1+p(p-1)\lambda^{p-2}
    \not\equiv 1\pmod{p^2},
\]
because $p\nmid \lambda$. Hence $\lambda+p$ has order $p(p-1)$ modulo $p^2$.

Therefore there exists an integer $r$ that is a primitive root modulo $p^2$.

Now we show that the same $r$ is a primitive root modulo $p^e$ for every $e\geq2$.

Since $r$ has order $p(p-1)$ modulo $p^2$, we have
\[
    r^{p-1}\equiv 1+cp\pmod{p^2}
\]
for some integer $c$ with $p\nmid c$.

We claim that for every $k\geq2$,
\[
    r^{p^{k-2}(p-1)}
    \equiv 1+c_kp^{k-1}\pmod{p^k},
    \qquad p\nmid c_k.
\]
The case $k=2$ is exactly the congruence above.

Assume the congruence holds for some $k\geq2$. Then
\[
    r^{p^{k-1}(p-1)}
    =
    \left(r^{p^{k-2}(p-1)}\right)^p
    \equiv
    (1+c_kp^{k-1})^p
    \pmod{p^{k+1}}.
\]
By the binomial theorem,
\[
    (1+c_kp^{k-1})^p
    \equiv
    1+c_kp^k
    \pmod{p^{k+1}},
\]
because the higher-order terms are divisible by $p^{k+1}$. Since $p\nmid c_k$, the induction continues.

Thus, modulo $p^e$,
\[
    r^{p^{e-2}(p-1)}\not\equiv 1\pmod{p^e},
\]
while Euler's theorem gives
\[
    r^{p^{e-1}(p-1)}\equiv1\pmod{p^e}.
\]
Therefore
\[
    \ord_{p^e}(r)=p^{e-1}(p-1).
\]
Hence $\Unit(p^e)$ is cyclic of order $p^{e-1}(p-1)$.
\end{proof}

\begin{proposition}[Alternative proof using Sylow subgroups]\label{prop:odd-prime-power-sylow}
Let $p$ be an odd prime and let $d\geq1$. Then $\Unit(p^d)$ is cyclic.
\end{proposition}

\begin{proof}
We use the structure theorem for finite abelian groups.

First, the Sylow $p$-subgroup of $\Unit(p^d)$ is cyclic. Indeed, one checks that $1+p$ has order $p^{d-1}$ modulo $p^d$. More precisely,
\[
    (1+p)^{p^{d-1}}\equiv1\pmod{p^d},
\]
but
\[
    (1+p)^{p^{d-2}}\not\equiv1\pmod{p^d}
\]
when $d\geq2$. Hence the Sylow $p$-subgroup is cyclic.

Now let $q\neq p$ be a prime. Consider the natural reduction map
\[
    \Z/p^d\Z\longrightarrow \Z/p\Z,
    \qquad
    a+(p^d)\longmapsto a+(p).
\]
It induces a surjective group homomorphism
\[
    \Unit(p^d)\longrightarrow \Unit(p).
\]
The kernel has order $p^{d-1}$. Therefore, for $q\neq p$, every Sylow $q$-subgroup of $\Unit(p^d)$ injects into $\Unit(p)$.

But
\[
    \Unit(p)\cong C_{p-1}
\]
is cyclic. Hence every Sylow $q$-subgroup of $\Unit(p^d)$ is cyclic.

Thus all Sylow subgroups of the finite abelian group $\Unit(p^d)$ are cyclic. Since a finite abelian group is the direct product of its Sylow subgroups, and cyclic groups of pairwise coprime orders have cyclic direct product, $\Unit(p^d)$ is cyclic.
\end{proof}

\section{Powers of Two}

\begin{theorem}[Units modulo powers of two]\label{thm:two-power}
For powers of $2$, the unit groups are:
\[
    \Unit(2)=\{1\},
    \qquad
    \Unit(4)\cong C_2,
\]
and for $e\geq3$,
\[
    \Unit(2^e)\cong C_{2^{e-2}}\times C_2.
\]
Consequently,
\[
    \Aut(C_{2^e})\cong C_{2^{e-2}}\times C_2
    \qquad(e\geq3).
\]
\end{theorem}

\begin{proof}
The cases $e=1$ and $e=2$ are immediate:
\[
    \Unit(2)=\{1\},
    \qquad
    \Unit(4)=\{\overline{1},\overline{3}\}\cong C_2.
\]

Now assume $e\geq3$.

We first prove that $5$ has order $2^{e-2}$ in $\Unit(2^e)$. Since $5=1+4$, we claim that for every $m\geq0$,
\[
    5^{2^m}\equiv 1+2^{m+2}\pmod{2^{m+3}}.
\]
For $m=0$, this says
\[
    5\equiv1+4\pmod8,
\]
which is true. If the congruence holds for $m$, then
\[
    5^{2^{m+1}}
    =
    \left(5^{2^m}\right)^2
    \equiv
    (1+2^{m+2})^2
    \pmod{2^{m+4}}.
\]
Expanding gives
\[
    (1+2^{m+2})^2
    =
    1+2^{m+3}+2^{2m+4}
    \equiv
    1+2^{m+3}
    \pmod{2^{m+4}}.
\]
So the claim follows by induction.

Taking $m=e-2$ gives
\[
    5^{2^{e-2}}\equiv1\pmod{2^e},
\]
while taking $m=e-3$ gives
\[
    5^{2^{e-3}}\equiv1+2^{e-1}\not\equiv1\pmod{2^e}.
\]
Therefore
\[
    \ord_{2^e}(5)=2^{e-2}.
\]

Now define
\[
    K=\{\overline{a}\in\Unit(2^e):a\equiv1\pmod4\}.
\]
Exactly half of the odd residues modulo $2^e$ are congruent to $1$ modulo $4$, so
\[
    |K|=2^{e-2}.
\]
Since $5\equiv1\pmod4$, all powers of $5$ lie in $K$. But $\gen{5}$ has order $2^{e-2}$, so
\[
    K=\gen{5}.
\]

Every odd residue is congruent either to $1$ or to $3$ modulo $4$. Thus every unit modulo $2^e$ lies either in $K$ or in $-K$. Hence
\[
    \Unit(2^e)=K\gen{-1}.
\]
Also,
\[
    K\cap\gen{-1}=\{1\},
\]
because $-1\equiv3\pmod4$ and hence $-1\notin K$. Therefore
\[
    \Unit(2^e)=\gen{5}\times\gen{-1}
    \cong C_{2^{e-2}}\times C_2.
\]
\end{proof}

\begin{remark}[Three involutions in $\Aut(C_{2^e})$]
For $e\geq3$, the group
\[
    \Aut(C_{2^e})\cong \Unit(2^e)
\]
has exactly three elements of order $2$. If $C_{2^e}=\gen{g}$, these correspond to the following three units modulo $2^e$:
\[
    -1,\qquad
    1+2^{e-1},\qquad
    2^{e-1}-1.
\]
Thus the three automorphisms of order $2$ are
\[
    g\longmapsto g^{-1},
    \qquad
    g\longmapsto g^{1+2^{e-1}},
    \qquad
    g\longmapsto g^{2^{e-1}-1}.
\]
\end{remark}

\section{Chinese Remainder Theorem}

\subsection{Ring version}

\begin{theorem}[Chinese Remainder Theorem]\label{thm:crt-ring}
Let $R$ be a commutative ring with identity, and let
\[
    A_1,A_2,\ldots,A_k
\]
be ideals of $R$. Define
\[
    \Phi:R\longrightarrow R/A_1\times R/A_2\times\cdots\times R/A_k
\]
by
\[
    \Phi(r)=(r+A_1,r+A_2,\ldots,r+A_k).
\]
Then $\Phi$ is a ring homomorphism and
\[
    \Ker\Phi=A_1\cap A_2\cap\cdots\cap A_k.
\]

If the ideals are pairwise comaximal, meaning
\[
    A_i+A_j=R
    \qquad(i\neq j),
\]
then $\Phi$ is surjective and
\[
    A_1\cap A_2\cap\cdots\cap A_k=A_1A_2\cdots A_k.
\]
Consequently,
\[
    R/(A_1A_2\cdots A_k)
    \cong
    \prod_{i=1}^k R/A_i.
\]
\end{theorem}

\begin{proof}
It is clear that $\Phi$ is a ring homomorphism. Also,
\[
    r\in\Ker\Phi
    \iff
    r\in A_i\text{ for every }i
    \iff
    r\in A_1\cap\cdots\cap A_k.
\]

Now assume that the ideals are pairwise comaximal.

We first prove surjectivity. Fix $i$. Since $A_i$ is comaximal with every $A_j$ for $j\neq i$, repeated use of the identity
\[
    (I+J=R,\ I+K=R)\Longrightarrow I+JK=R
\]
shows that
\[
    A_i+\prod_{j\neq i}A_j=R.
\]
Thus we may choose
\[
    u_i\in A_i,
    \qquad
    e_i\in\prod_{j\neq i}A_j
\]
such that
\[
    u_i+e_i=1.
\]
Then
\[
    e_i\equiv1\pmod{A_i},
    \qquad
    e_i\equiv0\pmod{A_j}\quad(j\neq i).
\]

Given an arbitrary element
\[
    (r_1+A_1,\ldots,r_k+A_k)
    \in
    \prod_{i=1}^k R/A_i,
\]
set
\[
    r=r_1e_1+r_2e_2+\cdots+r_ke_k.
\]
Then
\[
    \Phi(r)=(r_1+A_1,\ldots,r_k+A_k),
\]
so $\Phi$ is surjective.

It remains to prove
\[
    A_1\cap\cdots\cap A_k=A_1\cdots A_k.
\]
The inclusion
\[
    A_1A_2\cdots A_k\subseteq A_1\cap\cdots\cap A_k
\]
is immediate.

For the reverse inclusion, it suffices by induction to prove the case of two comaximal ideals. Suppose $A+B=R$. Choose $a\in A$ and $b\in B$ such that
\[
    a+b=1.
\]
If $x\in A\cap B$, then
\[
    x=x(a+b)=xa+xb.
\]
Since $x\in B$ and $a\in A$, we have $xa\in AB$. Since $x\in A$ and $b\in B$, we have $xb\in AB$. Therefore $x\in AB$, so
\[
    A\cap B=AB.
\]
The general case follows by induction.
\end{proof}

\subsection{Integer version}

\begin{corollary}[Chinese Remainder Theorem for integers]\label{cor:crt-int}
If $m_1,m_2,\ldots,m_k$ are pairwise coprime positive integers, then
\[
    \Z/(m_1m_2\cdots m_k)\Z
    \cong
    \Z/m_1\Z\times\Z/m_2\Z\times\cdots\times\Z/m_k\Z.
\]
\end{corollary}

\begin{proof}
Apply Theorem \ref{thm:crt-ring} to $R=\Z$ and
\[
    A_i=m_i\Z.
\]
Pairwise coprimality of the $m_i$ is exactly the condition
\[
    m_i\Z+m_j\Z=\Z
    \qquad(i\neq j).
\]
\end{proof}

\begin{corollary}[Unit group version]\label{cor:crt-units}
If $\gcd(m,n)=1$, then
\[
    (\Z/mn\Z)^\times
    \cong
    (\Z/m\Z)^\times\times(\Z/n\Z)^\times.
\]
More generally, if $m_1,\ldots,m_k$ are pairwise coprime, then
\[
    \Unit(m_1m_2\cdots m_k)
    \cong
    \prod_{i=1}^k \Unit(m_i).
\]
\end{corollary}

\begin{proof}
The Chinese Remainder Theorem gives a ring isomorphism
\[
    \Z/(m_1m_2\cdots m_k)\Z
    \cong
    \prod_{i=1}^k \Z/m_i\Z.
\]
A ring isomorphism sends units to units and induces an isomorphism between the corresponding unit groups. Therefore
\[
    \Unit(m_1m_2\cdots m_k)
    \cong
    \prod_{i=1}^k \Unit(m_i).
\]
\end{proof}

\section{The General Case}

\begin{theorem}[General structure of $\Aut(C_n)$]\label{thm:general-aut-cn}
Let
\[
    n=p_1^{e_1}p_2^{e_2}\cdots p_t^{e_t}
\]
be the prime factorization of $n$. Then
\[
    \Aut(C_n)
    \cong
    \Unit(n)
    \cong
    \Unit(p_1^{e_1})\times\Unit(p_2^{e_2})\times\cdots\times\Unit(p_t^{e_t}).
\]
Moreover,
\[
    \Unit(p^e)\cong
    \begin{cases}
        C_{p^{e-1}(p-1)}, & p\text{ odd},\\[4pt]
        \{1\}, & p=2,\ e=1,\\[4pt]
        C_2, & p=2,\ e=2,\\[4pt]
        C_{2^{e-2}}\times C_2, & p=2,\ e\geq3.
    \end{cases}
\]
\end{theorem}

\begin{proof}
By Theorem \ref{thm:aut-cyclic},
\[
    \Aut(C_n)\cong\Unit(n).
\]
By the Chinese Remainder Theorem for unit groups,
\[
    \Unit(n)
    \cong
    \prod_{i=1}^t \Unit(p_i^{e_i}).
\]
The structure of each odd prime-power factor is given by Theorem \ref{thm:odd-prime-power}, and the structure of each power-of-two factor is given by Theorem \ref{thm:two-power}.
\end{proof}

\begin{remark}[Summary table]
\[
\begin{array}{c|c|c}
\toprule
\text{Cyclic group} & \text{Automorphism group} & \text{Order}\\
\midrule
C_n
&
\Unit(n)
&
\varphi(n)
\\[3pt]
C_p
&
C_{p-1}
&
p-1
\\[3pt]
C_{p^e},\ p\text{ odd}
&
C_{p^{e-1}(p-1)}
&
p^{e-1}(p-1)
\\[3pt]
C_{2^e},\ e\geq3
&
C_{2^{e-2}}\times C_2
&
2^{e-1}
\\
\bottomrule
\end{array}
\]
\end{remark}

